% 2014 年江苏卷解答题第 16 题：三棱锥 P-ABC，D、E、F 分别为 PC、AC、AB 中点。
% PA⊥AC，PA=6，BC=8，DF=5。来源 PDF：普通高考/2014/2014江苏.pdf，印面 889。
\begin{tikzpicture}[>=Stealth, semithick, font=\small, scale=0.95]
  \coordinate (P) at (0,4.6);
  \coordinate (A) at (0,1.5);
  \coordinate (B) at (1.3,0);
  \coordinate (C) at (4.6,1.7);
  \coordinate (D) at ($(P)!0.5!(C)$);
  \coordinate (E) at ($(A)!0.5!(C)$);
  \coordinate (F) at ($(A)!0.5!(B)$);
  % 棱
  \draw (P) -- (A);
  \draw (P) -- (B);
  \draw (P) -- (C);
  \draw (A) -- (B);
  \draw (B) -- (C);
  \draw[dashed] (A) -- (C);
  % 辅助线
  \draw[dashed] (D) -- (E);
  \draw[dashed] (E) -- (F);
  \draw[dashed] (F) -- (D);
  \draw[dashed] (B) -- (E);
  \draw[dashed] (B) -- (D);
  % 顶点与中点标注
  \fill (P) circle (1.1pt) node[above] {$P$};
  \fill (A) circle (1.1pt) node[left] {$A$};
  \fill (B) circle (1.1pt) node[below] {$B$};
  \fill (C) circle (1.1pt) node[right] {$C$};
  \fill (D) circle (1.1pt) node[above right] {$D$};
  \fill (E) circle (1.1pt) node[below right] {$E$};
  \fill (F) circle (1.1pt) node[below left] {$F$};
\end{tikzpicture}
